% Fichier généré automatiquement par tools/generate_corrections.py.
% Source : SLCI/SLCI-03-SchemaBlocs/51_MCC/51_MCC.tex
% Statut : complet (2/2 question(s) renseignée(s), 0 reprise(s)).
\begin{corrigebox}[Corrigé extrait de la banque]
\CorrigeQuestion{1}
\begin{marginfigure}
\centering
\includegraphics[width=\linewidth]{51_01_c}

\end{marginfigure}
\CorrigeQuestion{2}
\begin{marginfigure}
\centering
\includegraphics[width=\linewidth]{51_01}

\end{marginfigure}

En utilisant le schéma-blocs proposé, on a $\Omega(p) = \left(C_r(p)A(p)+U(p)B(p)\right)C(p)$.

D'autre part,  $\Omega(p)=\left( C_r(p) + \dfrac{K}{R+Lp}\left(U(p)-K\Omega(p) \right) \right) \dfrac{1}{f+Jp}$.

On a donc $\left(f+Jp \right)\Omega(p)= C_r(p) + U(p)\dfrac{K}{R+Lp}  $

$\Leftrightarrow \left(f+Jp \right)\Omega(p)+ \dfrac{K^2}{R+Lp}\Omega(p)  = C_r(p) + U(p)\dfrac{K}{R+Lp} $

$\Leftrightarrow \left(\left(f+Jp \right)+ \dfrac{K^2}{R+Lp}\right)\Omega(p)  = C_r(p) + U(p)\dfrac{K}{R+Lp} $

$\Leftrightarrow \dfrac{K^2+\left(f+Jp \right)\left(R+Lp \right)}{R+Lp}\Omega(p)  = C_r(p) + U(p)\dfrac{K}{R+Lp} $

$\Leftrightarrow \Omega(p)  = \left(C_r(p) + U(p)\dfrac{K}{R+Lp}\right)\dfrac{R+Lp}{K^2+\left(f+Jp \right)\left(R+Lp \right)} $.

Dés lors plusieurs schéma-blocs peuvent répondre à la question. Par exemple, $A(p)=1$, $B(p)=\dfrac{K}{R+Lp}$, $C(p)=\dfrac{R+Lp}{K^2+\left(f+Jp \right)\left(R+Lp \right)}$.

En poursuivant, on a aussi : 
$ \Omega(p)  = \left(C_r(p)(R+Lp) + U(p) K\right)\dfrac{1}{K^2+\left(f+Jp \right)\left(R+Lp \right)} $.

On a donc aussi,  $A(p)=R+Lp$, $B(p)={K}$, $C(p)=\dfrac{1}{K^2+\left(f+Jp \right)\left(R+Lp \right)}$

\begin{enumerate}
    \item .
    \item $A(p)=R+Lp$, $B(p)={K}$, $C(p)=\dfrac{1}{K^2+\left(f+Jp \right)\left(R+Lp \right)}$ (plusieurs réponses possibles). 
\end{enumerate}
\end{corrigebox}
